
The first line contains two space-separated integers $N$ $(2 \le N \le 10^5)$ and $K$ $(1 \le K \le N - 1)$: the number of buried gemstones (and oscilloscope readings) and number of those gemstones that have already been discovered by other rockhounds, respectively.

Each of the next $K$ lines contains three space-separated integers $x_i$ $(|x_i| \le 10^6)$, $y_i$ $(-10^6 \le y_i < 0)$, and $f_i$ $(1 \le f_i \le N)$, describing the location and frequency of a gemstone that has already been discovered.  It is possible that two or more gemstones occupy the same location. It is guaranteed that the $f_i$ values are unique.

Finally, each of the last $N$ lines contains two space-separated integers $x_j$ $(|x_j| \le 10^6)$ and $f_j$ $(1 \le f_j \le N)$, describing an oscilloscope reading. It is guaranteed that the $x_j$ values are unique, that they are listed in increasing order, and that every gemstone (discovered and undiscovered) has exactly one oscilloscope reading. It is also guaranteed that there is at least one plausible configuration of gemstones.

